\documentclass[aspectratio=169, 10pt]{beamer}
\usetheme{metropolis}

% --- Edinburgh colours ----------------------------------------
\definecolor{UoERed}{HTML}{7A2318}
\definecolor{UoEGold}{HTML}{B8860B}
\definecolor{CodeBG}{HTML}{F7F3ED}
\definecolor{CodeFrame}{HTML}{D5C9B1}

\setbeamercolor{palette primary}{bg=UoERed, fg=white}
\setbeamercolor{frametitle}{bg=UoERed, fg=white}
\setbeamercolor{title separator}{fg=UoEGold}
\setbeamercolor{progress bar}{fg=UoEGold, bg=UoEGold!20}
\setbeamercolor{alerted text}{fg=UoERed}
\setbeamercolor{block title}{bg=UoERed!15, fg=UoERed}
\setbeamercolor{block body}{bg=UoERed!5}

% --- Packages -------------------------------------------------
\usepackage[T1]{fontenc}
\usepackage{lmodern}
\usepackage{amsmath, amssymb, mathtools}
\usepackage{listings}
\usepackage{graphicx, booktabs}
\usepackage{tikz}
\usetikzlibrary{arrows.meta, positioning, calc, decorations.pathreplacing}
\usepackage{hyperref}
\hypersetup{colorlinks=true, linkcolor=UoERed, urlcolor=UoERed!70!black}
\setbeamertemplate{navigation symbols}{}

\lstset{
  language=Python,
  basicstyle=\ttfamily\scriptsize,
  keywordstyle=\color{UoERed}\bfseries,
  commentstyle=\color{gray}\itshape,
  stringstyle=\color{UoEGold},
  backgroundcolor=\color{CodeBG},
  frame=single,
  rulecolor=\color{CodeFrame},
  numbers=none, breaklines=true, showstringspaces=false, tabsize=4,
  xleftmargin=4pt, xrightmargin=4pt, aboveskip=6pt, belowskip=4pt,
}

\AtBeginSection[]{
  \begin{frame}
    \vfill\centering
    \usebeamerfont{title}\insertsectionhead\par
    \vfill
  \end{frame}
}

\title{Week 3 --- Automatic Differentiation for Economics}
\subtitle{Forward Mode, Reverse Mode \& PyTorch Autograd}
\author{Dr Juan Zurita}
\institute{Deep Learning for Macroeconomics\\
The University of Edinburgh~$\cdot$~School of Economics}
\date{}

\begin{document}
\maketitle

\begin{frame}{Why Differentiation Matters}
Economic models are systems of equations involving derivatives:

\bigskip
\begin{itemize}
\item Euler equations: $u'(c_t) = \beta\,\mathbb{E}_t[u'(c_{t+1})(1 + r_{t+1})]$
\item FOCs of optimisation problems
\item Jacobians for stability analysis
\end{itemize}

\bigskip
\alert{Automatic differentiation} computes exact derivatives --- not numerical approximation, not symbolic algebra.
\end{frame}

\begin{frame}{Three Ways to Differentiate}
\begin{center}
\begin{tabular}{lll}
\toprule
Method & Pros & Cons \\
\midrule
Symbolic & Exact, interpretable & Expression swell \\
Numerical & Simple (finite diff.) & Approximation error \\
Automatic & Exact, efficient & Requires code graph \\
\bottomrule
\end{tabular}
\end{center}

\bigskip
\textbf{Autodiff} decomposes any program into elementary operations and applies the chain rule to each.
\end{frame}

\begin{frame}{Forward vs Reverse Mode}
\textbf{Forward mode}: propagate derivatives input $\to$ output.

Cost: one pass per input variable. Good when $n_{\text{inputs}} \ll n_{\text{outputs}}$.

\bigskip
\textbf{Reverse mode} (backpropagation): propagate derivatives output $\to$ input.

Cost: one pass per output variable. Good when $n_{\text{outputs}} \ll n_{\text{inputs}}$.

\bigskip
\begin{block}{Key insight}
Neural networks have \textbf{one} scalar loss and \textbf{millions} of parameters $\Rightarrow$ reverse mode is vastly more efficient.
\end{block}
\end{frame}

\begin{frame}{PyTorch Autograd}
\begin{block}{In practice}
\texttt{torch.autograd} builds a computational graph and computes gradients automatically.
\end{block}

\bigskip
Key functions:
\begin{itemize}
\item \texttt{torch.autograd.grad()} --- compute $\partial y / \partial x$
\item \texttt{torch.autograd.functional.jacobian()} --- full Jacobian matrix
\item \texttt{torch.autograd.functional.hessian()} --- second derivatives
\end{itemize}

\bigskip
This is what makes DEQNs possible: we can differentiate through the entire economic model.
\end{frame}

\begin{frame}{Summary}
\begin{itemize}
\item Autodiff gives \textbf{exact} derivatives at machine precision
\item Reverse mode is efficient for scalar-output functions (losses)
\item PyTorch autograd handles everything behind the scenes
\item We can differentiate Euler equations, value functions, and equilibrium conditions
\end{itemize}

\bigskip
\textbf{Next week:} Deep Equilibrium Networks --- using all of this to solve dynamic models.
\end{frame}

\end{document}
